\section{Memory scores, decision policy and adapter bounds}
\label{app:methodref}

\subsection{What enters the memory}

The memory accepts an embedding and an observed outcome $o\in[-1,1]$.
The full wrapper starts from a supplied backbone vector, applies a two-layer
ReLU projection and a chain of residual adapters, and stores the resulting
normalized vector. The projection starts with random weights. Its optional
offline training objective is a contrastive softmax loss on same-polarity
partners and opposite-polarity negatives. The direct memory interface accepts
a representation chosen by the caller. Neither entry point trains the backbone.

A record is keyed by case identifier, tenant and domain. The wrapper hashes
the supplied context to obtain the identifier. On a repeated key with $n$
previous observations, the memory increments the confirmation count, refreshes
the timestamp and replaces the outcome with $(n o_{\rm old}+o_{\rm new})/(n+1)$.
It also normalizes $n z_{\rm old}+z_{\rm new}$ after averaging. Confirmations
count submissions under a key and need not come from independent observers.
An outcome below $-0.3$ contributes failure evidence, and one above $0.3$
contributes success evidence. Every stored record contributes to familiarity.

\subsection{Evidence scores for one memory}

Let $s_i$ be the dot product of the normalized query and record $i$.
The implementation divides by the vector norm plus $10^{-10}$, so a zero
vector has zero similarity. Let $d_i=\exp(-\lambda a_i)$, where $a_i$ is
the integer age in days, and let $n_i$ be the confirmation count.
The severity $v_i$ is $5,4,3,2,1$ on the successive outcome intervals
$[-1,-0.8)$, $[-0.8,-0.5)$, $[-0.5,0)$, $[0,0.5)$ and $[0.5,1]$.

From failures with $s_i>\tau$, select up to 20 in decreasing similarity
order to form $F$. Write $[x]_a^b=\min(b,\max(a,x))$. The risk score is
\begin{align}
 w_i &= [\alpha\,s_i(v_i/5)d_i n_i]_0^{0.99},\\
 b_F &=
 \begin{cases}
  0.1\max_{i<j,\ i,j\in F} z_i^\top z_j & |F|\geq 2,\\
  0 & |F|<2,
 \end{cases}\\
 R &= \left[1-\prod_{i\in F}(1-w_i)+b_F\right]_0^1 .
\end{align}
For all successes $S$ with $s_i>\tau$, the success score is
\[
 U=\left[\frac{1}{|S|}\sum_{i\in S}s_i\,\frac{6-v_i}{5}\,d_i\right]_0^1 .
\]
An empty evidence set gives zero for its score. The two scores can both be
high. Shared neighbors, repeated confirmations and the pairwise term make
a probability interpretation of the noisy-OR unwarranted.

\subsection{Fixed-bandwidth familiarity}

For a memory with $N\geq 10$ records, the default coverage score uses bandwidth
$h=0.3$. Define
\[
 K_h(z,x)=\exp\!\left(-\frac{(1-z^\top x)^2}{2h^2}\right),\qquad
 D_N(z)=\frac{1}{N}\sum_{i=1}^N K_h(z,z_i).
\]
The memory evaluates $D_N$ at up to 256 stored vectors selected at evenly
spaced record indices, including the selected records' own contributions.
If $q$ is the 90th percentile of those self-scores, coverage is
$C(z)=[D_N(z)/\max(q,10^{-12})]_0^1$. Memories with fewer than ten
records return zero coverage. Rescaling is per memory and clipping can tie
queries with different unscaled scores.

For each fixed query $z$ and fixed $h>0$, if the stored $Z_i$ are independent
draws from a common distribution $P$ on the unit sphere, the bounded summands
obey the strong law of large numbers. Thus
\[
 D_N(z)\ \longrightarrow\ \int K_h(z,x)\,dP(x)
 \qquad\text{almost surely}.
\]
The limit measures kernel-weighted probability mass. For example, if
$P$ is a point mass at $u$, a query orthogonal to $u$ receives the positive
limit $\exp[-1/(2h^2)]$ despite its neighborhood having zero mass.
Identifying a pointwise density would require further assumptions, bandwidth
selection and an appropriate normalization. Deduplication, adaptive feedback
and pruning can also violate the sampling assumption. The result concerns the
unscaled kernel mean. It establishes neither convergence of the capped-sample
calibration quantile nor discrimination between arbitrary distributions.

\subsection{Aggregation, indexing and actions}

The manager reads the global memory, the requested domain memory if present,
and the requested tenant memory if present. It takes the maximum risk, maximum
success and arithmetic mean coverage across these memories, including an empty
global memory. The caller supplies trusted scope identifiers. This partitioning
rule assumes the application has already checked access to those scopes.

With a cached $N$-by-$d$ embedding matrix, a dense similarity pass takes
$O(Nd)$ arithmetic. Coverage also averages $N$ kernel values. Failure retrieval
can sort as many as $N$ candidates and computes the pairwise term on at most
20 vectors. Recalibrating after a write requires up to $\min(N,256)$ extra
similarity passes. The optional scikit-learn cosine neighbor index uses brute
force for unfiltered top-$k$ retrieval. Outcome-filtered retrieval and the
three decision scores use the dense matrix. There is no logarithmic query-time
guarantee for the released implementation.

Table~\ref{tab:release-policy} gives the first matching rule for the default
consensus engine. The user supplies deterministic heuristic checks. At fixed
risk, success and heuristic results, raising coverage can take a case out of
the Query region once. It cannot return the case to Query at higher coverage.
The numeric confidence attached to an action is a policy output.
\begin{table}[htbp]
\centering
\begin{tabular}{ll}
\toprule
Condition in priority order & Action\\
\midrule
Critical rule violated or failed to evaluate & Escalate\\
$C<0.2$ & Query\\
$R>0.85$ & Escalate\\
$R>0.6$ & Abstain\\
Noncritical rule violated & Abstain\\
$U>0.7$ & Trust\\
$C<0.5$ & Query for validation\\
Otherwise & Trust at low confidence\\
\bottomrule
\end{tabular}
\caption{Action thresholds and precedence. Equalities follow the next applicable rule.}
\label{tab:release-policy}
\end{table}

\subsection{What the adapter update bounds}

Let $\theta$ concatenate the trainable parameters of one residual adapter,
including the affine LayerNorm parameters and biases. For an embedding $z$,
the adapter normalizes $z+W_2\operatorname{ReLU}(W_1\operatorname{LN}(z)
+b_1)+b_2$. With $s$ the dot product of that output with the normalized replay
input, the example loss is $1-s$ for $o>0$ and $\max(0,s+0.5)$ otherwise.
The batch objective adds $\lambda_{\rm reg}\sum_j\|\theta^{(j)}\|_2$,
where each parameter tensor is flattened separately. The adapter clips
the gradient of the entire objective jointly, then applies plain SGD with
zero momentum and zero optimizer weight decay. The code reports the mean
example loss before regularization.

Let $B=\{v:\|v\|_2\leq r\}$ and $P_B$ be its Euclidean projection.
Initialization projects all parameters jointly before recording the baseline
or copying the exponential moving average. With a fixed radius, a finite
clipped gradient $\|g_t\|_2\leq c$ and step size $\eta_t\geq 0$, the update is
\[
 \theta_0=P_B(\theta_{\rm init}),\qquad
 \theta_{t+1}=P_B(\theta_t-\eta_t g_t).
\]
Since $P_B(\theta_t)=\theta_t$ and projection onto a closed convex ball
is nonexpansive,
\[
 \|\theta_{t+1}-\theta_t\|_2
 \leq\|(\theta_t-\eta_t g_t)-\theta_t\|_2\leq\eta_t c,\qquad
 \sum_{t=0}^{T-1}\|\theta_{t+1}-\theta_t\|_2\leq c\sum_{t=0}^{T-1}\eta_t .
\]
These bounds hold in exact arithmetic. The measured path length
aggregates the post-projection steps and excludes initialization. Endpoint
displacement is also bounded by $2r$. A raw weight import starts a new
measurement interval at its projected state. Loading an older adapter checkpoint
likewise establishes a fresh baseline for subsequent updates.

The moving average starts in $B$ and remains there by convexity for
$0\leq\beta\leq1$. Direct adapter calls use live weights by default and the
wrapper's inference chain uses the moving average. An output-displacement
bound at a fixed input would additionally require a Lipschitz bound in
\emph{parameters} along the relevant path. The parameter bound alone
establishes no decision agreement, task accuracy or retention guarantee.
